\documentclass{article}
\usepackage{amsmath}
\usepackage{amsfonts}
\usepackage{amssymb}
\usepackage{graphicx}
\usepackage{hyperref}
\usepackage[a4paper, margin=1.0in]{geometry}


\title{Singular Perturbations, Moments and Descriptor Systems: Linear}
\author{Yoonjae Hwang}
\date{\today}

\begin{document}

\maketitle

\section{Linear Systems - Thinking about explicit solutions and moments}
Consider the following linear time-invariant, SISO system:
\label{sec:linear-systems}
\begin{equation}
\begin{aligned}
\dot{x} &= A x + B u, \\
y & = Cx,
\end{aligned}
\end{equation}

\begin{equation}
\begin{aligned}
\dot{w} &= Sw, \\
u &= Lw,
\end{aligned}
\end{equation}
Given initial conditions \(x(0) = x_0\) and \(w(0) = w_0\), the solution to the system is given by
\begin{equation}
x(t) = e^{At} x_0 + \int_0^t e^{A(t-\tau)} B L e^{S\tau} w_0 d\tau.
\end{equation}
The output of the system is given by
\begin{equation}
y(t) = C e^{At} x_0 + \int_0^t C e^{A(t-\tau)} B L e^{S\tau} w_0 d\tau.
\end{equation}
If $\sigma(A) \cup \sigma(S) = \emptyset$, then there exists a unique solution $\Pi$ to the Sylvester equation
\begin{equation}
A \Pi + B L = \Pi S.
\end{equation}
Using this, we can substitute an expression for \(B L\) into the integral term for $x(t)$:
\begin{equation}
\begin{aligned}
x(t) &= e^{At} x_0 + \int_0^t e^{A(t-\tau)} (\Pi S - A \Pi) e^{S\tau} w_0 d\tau \\
&= e^{At} x_0 + \int_0^t e^{A(t-\tau)} \Pi S e^{S\tau} w_0 d\tau - \int_0^t e^{A(t-\tau)} A \Pi e^{S\tau} w_0 d\tau \\
&= e^{At} x_0 + \left( \Pi e^{St} - e^{At} \Pi \right) w_0 .
\end{aligned}
\end{equation}
Forward Moments are defined to be 
\begin{equation}
M_g = C \Pi
\end{equation}

\section{Problem Setup}
\label{sec:problem-setup}
Consider the following linear singularly perturbed system:
\begin{equation}
\begin{aligned}
\dot{x}_1(t) &= A_{11} x_1(t) + A_{12} x_2(t) + B_1 u(t), \\
\varepsilon \dot{x}_2(t) &= A_{21} x_1(t) + A_{22} x_2(t) + B_2 u(t),
\end{aligned}
\end{equation}
where \(x_1(t) \in \mathbb{R}^{n_1}\), \(x_2(t) \in \mathbb{R}^{n_2}\), \(u(t) \in \mathbb{R}^m\), and \(0 < \epsilon \ll 1\).

You can write this as a descriptor system: 
\begin{equation}
    E_1 \dot{x} = A x + B u
\end{equation}
where
\begin{equation}
    E_1 = \begin{bmatrix}
    I_{n_1} & 0 \\
    0 & \varepsilon I_{n_2}
    \end{bmatrix}, \quad
    A = \begin{bmatrix}
    A_{11} & A_{12} \\
    A_{21} & A_{22}
    \end{bmatrix}, \quad
    B = \begin{bmatrix}
    B_1 \\
    B_2
    \end{bmatrix}, \quad
    x = \begin{bmatrix}
    x_1 \\
    x_2
    \end{bmatrix}.
\end{equation}
The input signal is also a linear singularly perturbed system: 
\begin{equation}
\begin{aligned}
\dot{w_1} = S_{11} w_1 + S_{12}w_2 \\
\varepsilon \dot{w_2} = S_{21} w_1 + S_{22} w_2
\end{aligned}
\end{equation}
where \(w_1(t) \in \mathbb{R}^{r_1}\), \(w_2(t) \in \mathbb{R}^{r_2}\), and \(0 < \epsilon \ll 1\). 
The input to the system is given by \(u(t) = L_1 w_1(t) + L_2 w_2(t)\).

The input signal can also be written as a descriptor system like so:
\begin{equation}
    E_2 \dot{w} = S w
\end{equation}
where
\begin{equation}
    E_2 = \begin{bmatrix}
    I_{r_1} & 0 \\
    0 & \varepsilon I_{r_2}
    \end{bmatrix}, \quad
    S = \begin{bmatrix}
    S_{11} & S_{12} \\
    S_{21} & S_{22}
    \end{bmatrix}, \quad
    w = \begin{bmatrix}
    w_1 \\
    w_2
    \end{bmatrix}.
\end{equation}
The output of the system is given by
\begin{equation}
    y(t) = C_1 x_1(t) + C_2 x_2(t).
\end{equation}

\section{Weierstrass Canonical Form}
\label{sec:weierstrass-canonical-form}
For a descriptor system of the form \(E \dot{x} = A x + B u\), the matrix pencil \((A, E)\)
is said to be regular
if there exists a complex number \(s\) such that \(\det(sE - A) \neq 0\).
If the matrix pencil \((A, E)\) is regular, then there exist nonsingular
matrices \(P\) and \(Q\) such that
\begin{equation}
P E Q = \begin{bmatrix}
I_{n_1} & 0 \\
0 & N
\end{bmatrix}, \quad
P A Q = \begin{bmatrix}
J & 0 \\
0 & I_{n_2}
\end{bmatrix},
\end{equation}
where \(J\) is in Jordan canonical form, and \(N\) is a nilpotent matrix. 

Let's assume $(A, E_1)$ and $(S, E_2)$ are regular. So for the singular perturbation system, we have 
\begin{equation}
P E_1 Q = \begin{bmatrix}
I_{n_1} & 0 \\
0 & N
\end{bmatrix}, \quad
P A Q = \begin{bmatrix}
J_A & 0 \\
0 & I_{n_2}
\end{bmatrix},
\end{equation}
and
\begin{equation}
W E_2 T = \begin{bmatrix}
I_{r_1} & 0 \\
0 & M
\end{bmatrix}, \quad
W S T = \begin{bmatrix}
J_S & 0 \\
0 & I_{r_2}
\end{bmatrix},
\end{equation}
where \(J_A\) and \(J_S\) are in Jordan canonical form, and \(N\) and \(M\) are nilpotent matrices. 
\newline
\textbf{Question}: Can we find P and Q for this? How does this relate to $A_{22}$ beign Hurwitz?

We can rewrite the system with a coordinate transformation \(x = Q \tilde{x}\) and \(w = T \tilde{w}\), 
with $\tilde{x} = \begin{bmatrix}\tilde{x}_1 \\ \tilde{x}_2 \end{bmatrix}$ and
$\tilde{w} = \begin{bmatrix}\tilde{w}_1 \\ \tilde{w}_2 \end{bmatrix}$:
\begin{equation}
\begin{aligned}
\begin{bmatrix}
I_{n_1} & 0 \\
0 & N
\end{bmatrix} \begin{bmatrix}
\dot{\tilde{x}}_1(t) \\
\dot{\tilde{x}}_2(t)
\end{bmatrix} &= \begin{bmatrix}
J_A & 0 \\
0 & I_{n_2}
\end{bmatrix} \begin{bmatrix}
\tilde{x}_1(t) \\
\tilde{x}_2(t)
\end{bmatrix} + \tilde{B} (L_1 \tilde{w}_1(t) + L_2 \tilde{w}_2(t)), \\
\begin{bmatrix}
I_{r_1} & 0 \\
0 & M
\end{bmatrix} \begin{bmatrix}
\dot{\tilde{w}}_1(t) \\
\dot{\tilde{w}}_2(t)
\end{bmatrix} &= \begin{bmatrix}
J_S & 0 \\
0 & I_{r_2}
\end{bmatrix} \begin{bmatrix}
\tilde{w}_1(t) \\
\tilde{w}_2(t)
\end{bmatrix}, \\
y(t) &= \tilde{C}_1 \tilde{x}_1(t) + \tilde{C}_2 \tilde{x}_2(t),
\end{aligned}
\end{equation}




\bibliographystyle{plain}
\bibliography{referenceses}



\end{document}